% Generated draft: requires mathematical and source review.
\documentclass[11pt]{article}
\usepackage[margin=25mm]{geometry}
\usepackage{fontspec}
\setmainfont{Latin Modern Roman}
\usepackage{amsmath,amssymb,mathtools,unicode-math}
\setmathfont{Latin Modern Math}
\usepackage{enumitem,needspace,xurl,hyperref}
\setlength{\parindent}{0pt}
\setlength{\parskip}{5pt}
\setlength{\emergencystretch}{3em}
\newcommand{\N}{\mathbb{N}}
\newcommand{\Z}{\mathbb{Z}}
\newcommand{\Q}{\mathbb{Q}}
\newcommand{\R}{\mathbb{R}}
\newcommand{\C}{\mathbb{C}}
\newcommand{\F}{\mathbb{F}}
\newcommand{\A}{\mathbb{A}}
\newcommand{\PP}{\mathbb{P}}
\newcommand{\E}{\mathbb{E}}
\newcommand{\eps}{\varepsilon}
\newcommand{\dd}{\,\mathrm{d}}
\begin{document}
\section*{UnsolvedMath ID: 30004853}
\textbf{Source problem code:} catalogue entry
\subsection{Definitions and mathematical statement}
Write \(\phi_f\) for the monodromy denoted in the supplied record, \(HF^*(\phi_f^d)\) for its indicated Floer cohomology, and \(\chi_d(f)\) for the indicated contact locus of jets. Here \(H_c^*\) denotes compactly supported cohomology. The record does not specify the domain or hypotheses on \(f\), the meaning of \(n\), the definition of \(\chi_d(f)\), the range of \(d\), the coefficients, or the Floer-theoretic conventions. These must be recovered from the cited source before this becomes a fully specified geometric problem.

Subject to those unresolved qualifications, let \(D\) denote the intended index set and put \(s_d=2nd+n-1\). Assume that compatible choices of the asserted graded isomorphisms have been made:
\[
I_d^k:HF^k(\phi_f^d)\xrightarrow{\ \cong\ }H_c^{k+s_d}(\chi_d(f)),\qquad d\in D.
\]
Define graded groups
\[
\mathcal F^k=\bigoplus_{d\in D}HF^k(\phi_f^d),\qquad
\mathcal H^k=\bigoplus_{d\in D}H_c^{k+s_d}(\chi_d(f)),
\]
and let \(I=\bigoplus_{d\in D}I_d\). Denote the pair-of-pants product on \(\mathcal F\) by \(\star\).

\textbf{Problem.} Identify, in terms of the geometry of the contact loci of jets, the operation \(\mu\) on \(\mathcal H\), or equivalently on the source's direct sum \(\bigoplus_d H_c^*(\chi_d(f))\) with these grading shifts understood, that satisfies
\[
\mu(Ix,Iy)=I(x\star y)
\]
for every pair of inputs for which the pair-of-pants product is defined. This identity specifies the desired compatibility; the question asks for a geometric description of the corresponding operation. The source excerpt does not fix its degree, its rule for output indices, or the choices needed to specify the isomorphisms, so no such conventions are imposed here.

\textbf{Qualifications from the supplied excerpt.} The isomorphism is stated there as Conjecture 2, attributed to [BdBLN20]. The excerpt reports that the same authors verified it when \(d\) equals the multiplicity of \(f\) at \(0\), and then announces general validity as ``Theorem 3 (M., in progress).'' These are historical source claims, not independently verified results in this draft. The proposed proof constructs a natural PSS morphism
\[
\operatorname{ev}:HF^*(\phi_f^d)\longrightarrow H_c^{*+2nd+n-1}(\chi_d(f)).
\]
The excerpt describes holomorphic disks whose boundaries limit to an orbit of the mapping torus of \(\phi_f^d\), with their \(d\)-jets sweeping out a cycle in a thickening of \(\chi_d(f)\) inside the jet space written there as \(\operatorname{Jet}_d(f)|_0\). It compares this construction with the logarithmic PSS map of Ganatra and Pomerleano [GP20]. It further proposes filtrations \(F,F'\), arising from spectral sequences in [McL19] and [BdBLN20], such that \(\operatorname{ev}\) preserves the filtrations and induces an isomorphism on associated graded objects. These proposed constructions do not supply a product correspondence in the provided excerpt.
\subsection{Short English statement}
Assuming suitable arc--Floer isomorphisms, describe geometrically the operation on compactly supported cohomology of jet contact loci that corresponds to the Floer pair-of-pants product, keeping the shifts \(2nd+n-1\). The supplied record leaves essential geometric hypotheses and conventions unspecified, and its general isomorphism claim is presented historically as work in progress.
\subsection{Sources}
\begin{itemize}
\item Local UnsolvedMath registry, supplied complete source record, permanent ID \(30004853\), source problem code \texttt{catalogue entry}, title ``Pair-of-Pants Products on Contact Loci of Jets.'' The record belongs to the Oberwolfach Reports Open Problems collection; its proposed year is 2021 and no proposer is named. Its labels \texttt{open} and \texttt{corrected\_verified}, including the supplied claim of manual semantic, typographic, and LaTeX review, are registry annotations rather than independent mathematical verification.
\item Supplied report citation: \emph{Singularities (2022)}, Oberwolfach Reports, pp. 2573--2639, DOI 10.4171/owr/2021/47. \url{https://doi.org/10.4171/owr/2021/47}. The supplied citation uses 2022, while the proposed year and DOI contain 2021; both are retained.
\item Supplied literature source, ``OWR report posing the product question'': \url{https://ems.press/content/serial-article-files/46924}.
\item Supplied literature source, ``Cohomology of contact loci'': \url{https://doi.org/10.4310/jdg/1649953456}.
\item Supplied literature source, ``Recent special-case arc--Floer results'': \url{https://arxiv.org/abs/2506.18805}.
\item The supplied excerpt attributes Conjecture 2 and its multiplicity case to [BdBLN20], the announced general theorem to ``M., in progress,'' a related logarithmic PSS map to Ganatra and Pomerleano [GP20], and the relevant spectral sequences to [McL19] and [BdBLN20]. Full bibliographic entries and the identity of ``M.'' are not supplied.
\item Literature assessment dated exactly \texttt{2026-08-22}: the supplied curation labels the problem open, describes the report's product question as a further direction, and states that subsequent arc--Floer literature establishes cohomological cases without identifying the product in general. This is supplied dated curation, not an independently verified literature assessment; its time-sensitive status and cited claims require independent checking.
\end{itemize}
\end{document}
